%-------------------------------------------------------------
% Simple CV template
% Compile with: pdflatex cv_template.tex  (or upload to Overleaf)
%-------------------------------------------------------------
\documentclass[10.5pt,a4paper]{article}

%---------- Packages ----------
\usepackage[utf8]{inputenc}
\usepackage[T1]{fontenc}
\usepackage{lmodern}
\usepackage[margin=1.75cm]{geometry}
\usepackage{xcolor}
\usepackage[most]{tcolorbox}
\usepackage{enumitem}
\usepackage{titlesec}
\usepackage[hidelinks]{hyperref}

%---------- Colours ----------
\definecolor{accent}{HTML}{FCA311}       % headings and rules
\definecolor{boxbg}{HTML}{14213D}        % skill box background
\definecolor{boxtext}{HTML}{E5E5E5}      % skill box textcolor
\definecolor{boxborder}{HTML}{14213D}    % skill box border
\definecolor{specialgrey}{HTML}{636363}  % color for grey elements

%---------- Page setup ----------
\pagestyle{empty}                      % no page numbers
\setlength{\parindent}{0pt}
\setlength{\parskip}{0pt}
\raggedright

%---------- Section headings ----------
\titleformat{\section}
  {\large\bfseries\color{accent}\scshape}
  {}{0pt}
  {}
  [\vspace{-0.6em}\color{accent}\rule{\linewidth}{0.6pt}]
\titlespacing*{\section}{0pt}{1.2em}{0.6em}

%---------- Custom commands ----------

% Skill box:  \skill{Python}
\newcommand{\skill}[1]{%
  \tcbox[on line, arc=4pt, boxrule=0.5pt,
         coltext=boxtext, colback=boxbg,
         colframe=boxborder, left=4pt,
         right=4pt, height=1.25em,
         valign=center, boxsep=0pt,
         nobeforeafter]{%
         \small\rule[-0.35em]{0pt}{1.5em}#1}%
  \hspace{0.3em}\allowbreak}

% Entry with title/date line and subtitle/location line:
% \entry{Title}{Date}{Organisation}{Location}
\newcommand{\entry}[5][]{%
  \textbf{#2} \hfill #3 \\
  \textit{#4} \hfill \textit{#5}
  \if\relax\detokenize{#1}\relax\else\\ {\small #1}\fi
  \par\vspace{0.2em}}

% Compact bullet list for details under an entry
\newlist{details}{itemize}{1}
\setlist[details]{label=\textbullet, leftmargin=1.4em, itemsep=1pt, topsep=2pt, parsep=0pt}

%==============================================================
\begin{document}

%---------- 1. Header ----------
\begin{center}
  {\Huge\bfseries Levente László POCSAY}\\[0.6em]
  \small
  +43 678 7831628 \quad|\quad
  \href{mailto:lpocsay@gmail.com}{lpocsay@gmail.com} \quad|\quad
  \href{https://www.linkedin.com/in/pleventel}{linkedin.com/in/pleventel} \quad|\quad
  \href{https://www.pleventel.github.io/CV}{pleventel.github.io/CV}
\end{center}

%---------- 2. Profile + skills ----------
\section{Profile}
I am an Artificial Intelligence Solutions student and Research Trainee, currently starting my third semester at the University of Applied Sciences Upper Austria,
while also working remotely as a Research Trainee on a time-series analysis project for the HUN-REN Wigner Research Centre for Physics,
which has been an ongoing project for almost three years.

\vspace{0.8em}
\skill{\texttt{Python}}
\skill{\texttt{pandas}}
\skill{\texttt{NumPy}}
\skill{\texttt{matplotlib}}
\skill{\texttt{pytorch}}
\skill{\texttt{Altair}}
\skill{\texttt{SQL}}
\skill{\texttt{C++}}
\skill{Git}
\skill{\LaTeX}
\skill{Artificial Intelligence}
\skill{Machine Learning}
\skill{Neural Networks}
\skill{Data Analysis}

%---------- 3. Education ----------
\section{Education}
\entry{Artificial Intelligence Solutions}{Sep 2025 -- present}
      {University of Applied Sciences Upper Austria}{Hagenberg i. Mkr., Austria}
\begin{details}
  \item Courses in Supervised Machine Learning, Symbolic AI, Unsupervised Machine Learning, Neural Networks and Deep Learning.
  \item Programming in Python and C++.
  \item \textcolor{specialgrey}{Future courses (in SS 2027): Computer Vision, Embedded AI, Generative AI \& Natural Language Processing.}
\end{details}

%---------- 4. Experience ----------
\section{Experience}
\entry[Faculty of Informatics, Communication and Media]{Tutor}{Oct 2026 -- present}
      {University of Applied Sciences Upper Austria}{Hagenberg i. Mkr., Austria}
Working as a Tutor / Teaching Assistant in class MAT1 -- Basics of Linear Algebra and Calculus in the degree program Artificial Intelligence Solutions BSc.
\begin{details}
  \item Providing in-class help to the lecturer by guiding students with their tasks.
  \item Coordinating task distribution with other tutors and lecturer.
  \item Checking and marking students' uploaded assignemnts in Moodle.
\end{details}
\entry[Department of Compuational Sciences]{Research Trainee}{Sep 2025 -- present}
      {HUN-REN Wigner Research Centre for Physics}{Budapest, Hungary}
\begin{details}
  \item Working on time series analysis anomaly detection with the ALT method.
  \begin{itemize}
    \item Creating our own artificial dataset for our anomaly detection;
    \item Review previous research and publication;
    \item Finding time series databases -- containing anomaly -- appropriate for the project's purporse.
  \end{itemize}
  \item Mentoring \& guiding high school students in their research project involving first time-series classification, then later on edge detection (Computer Vision).
\end{details}

\vspace{0.6em}

\entry[Department of Compuational Sciences]{Research Assistant}{Jul 2026 -- Sep 2026}
      {HUN-REN Wigner Research Centre for Physics}{Budapest, Hungary}
\begin{details}
  \item Research-oriented software developement: created Python package \texttt{artificial-dataset} to be used for creating time-series anomalious datasets.
  \item Practical implementation and experimentation with machine-learning approaches for anomaly detection, including Isolation Forest, Local Outlier Factor, and autoencoder-based anomaly detection.
  \item Experience in evaluating anomaly-detection models using quantitative performance metrics, comparing different approaches and interpreting their results in a systematic manner.
\end{details}


%---------- 5. Projects (optional, delete if not needed) ----------
% \section{Projects}
% \entry{Project Name}{2025}
%       {\href{https://github.com/yourname/project}{github.com/yourname/project}}{}
% \begin{details}
%   \item One or two lines on what it does and which technologies it uses
% \end{details}
% 
% %---------- 6. Certifications & awards (optional) ----------
% \section{Certifications \& Awards}
% \begin{details}
%   \item \textbf{Certificate or Award Name} -- Issuer, 2025
%   \item \textbf{Certificate or Award Name} -- Issuer, 2024
% \end{details}

%---------- 7. Languages ----------
\section{Languages}
\begin{tabular}{@{}p{3cm}p{5cm}@{}}
  Hungarian  & Native \\
  English & Fluent (C1) \\
  German  & Under-intermediate (B1) \\
\end{tabular}

\end{document}